% GENERATED FILE. Edit canonical lessons, then run npm run generate:books.
% canonical-source-hash: fnv1a64:d26b12e58de1832d
% canonical-lessons: AR-C148-masrur, AR-C148-wasi, AR-C148-amiq, AR-C148-mustadir, AR-C148-hadd

\chapter{Pleased, Spacious, Deep, Round, Sharp}
\label{ch:ar-pleased-spacious-deep-round-sharp}
\hlchaptermodality{148}

\begin{chapteropening}
\textbf{By the end of this chapter:} \emph{I can say pleased, spacious, deep, round and sharp.}

Complete the last lesson of chapter 148: I can say pleased, spacious, deep, round and sharp.
\end{chapteropening}

\section[\texorpdfstring{\ar{مسرور} (masrūr)}{مسرور (masrūr)}]{\ar{مسرور} (masrūr) — pleased}

\label{lesson:AR-C148-masrur}



\begin{quote}
Before the new one: say the Arabic for kind, pleasant, then the Arabic for angry.
\end{quote}

\subsection*{You'll want to know: \ar{مسرور}}
\textbf{\ar{مسرور}} — \emph{masrūr} — \textquotedblleft{}pleased\textquotedblright{}.

\begin{cousinweb}
\textbf{\ar{مسرور}} (\emph{masrūr}) is built on the native Arabic root \emph{s-r-r}, \textquotedblleft{}to gladden\textquotedblright{} (Lane). \emph{Masrūr} is the passive participle, \textquotedblleft{}gladdened, pleased\textquotedblright{}.
\end{cousinweb}

\begin{grammarlens}[title={What you've built}]
One more word for this chapter.
\end{grammarlens}

\subsection*{Your turn}
\begin{itemize}
\raggedright
  \item \emph{Say it:} \emph{masrūr}
  \item \emph{Say it:} \emph{masrūr}, once more
  \item \emph{Say it:} say \emph{masrūr}
  \item \emph{Recall:} say the Arabic for kind, pleasant, then the Arabic for angry, then say \emph{masrūr} again
\end{itemize}

\begin{tcolorbox}[breakable,colback=teal!4,colframe=teal!35!black,title={Before you move on}]
What does \ar{مسرور} mean? (\textquotedblleft{}Pleased\textquotedblright{}.) Say it once more.
\end{tcolorbox}
\section[\texorpdfstring{\ar{واسع} (wāsiʿ)}{واسع (wāsiʿ)}]{\ar{واسع} (wāsiʿ) — spacious, wide}

\label{lesson:AR-C148-wasi}



\begin{quote}
Before the new one: say the Arabic for angry, then the Arabic for pleased.
\end{quote}

\subsection*{You'll want to know: \ar{واسع}}
\textbf{\ar{واسع}} — \emph{wāsiʿ} — \textquotedblleft{}spacious, wide\textquotedblright{}.

\begin{cousinweb}
\textbf{\ar{واسع}} (\emph{wāsiʿ}) is built on the native Arabic root \emph{w-s-ʿ}, \textquotedblleft{}to be wide\textquotedblright{} (Lane). \emph{Wāsiʿ} is the active participle, \textquotedblleft{}wide\textquotedblright{}.
\end{cousinweb}

\begin{grammarlens}[title={What you've built}]
One more word for this chapter.
\end{grammarlens}

\subsection*{Your turn}
\begin{itemize}
\raggedright
  \item \emph{Say it:} \emph{wāsiʿ}
  \item \emph{Say it:} \emph{wāsiʿ}, once more
  \item \emph{Say it:} say \emph{wāsiʿ}
  \item \emph{Recall:} say the Arabic for angry, then the Arabic for pleased, then say \emph{wāsiʿ} again
\end{itemize}

\begin{tcolorbox}[breakable,colback=teal!4,colframe=teal!35!black,title={Before you move on}]
What does \ar{واسع} mean? (\textquotedblleft{}Spacious, wide\textquotedblright{}.) Say it once more.
\end{tcolorbox}
\section[\texorpdfstring{\ar{عميق} (ʿamīq)}{عميق (ʿamīq)}]{\ar{عميق} (ʿamīq) — deep}

\label{lesson:AR-C148-amiq}



\begin{quote}
Before the new one: say the Arabic for pleased, then the Arabic for spacious.

Say \textquotedblleft{}It is cold.\textquotedblright{} (\emph{Al-jaw bārid}.) And the formal word for weather? (\emph{Aṭ-ṭaqs}.)
\end{quote}

\subsection*{You'll want to know: \ar{عميق}}
\textbf{\ar{عميق}} — \emph{ʿamīq} — \textquotedblleft{}deep\textquotedblright{}.

\begin{cousinweb}
\textbf{\ar{عميق}} (\emph{ʿamīq}) is built on the native Arabic root \emph{ʿ-m-q}, \textquotedblleft{}to be deep\textquotedblright{} (Lane). \emph{ʿAmīq} is the \emph{faʿīl} adjective, \textquotedblleft{}deep\textquotedblright{}.
\end{cousinweb}

\begin{grammarlens}[title={What you've built}]
One more word for this chapter.
\end{grammarlens}

\subsection*{Your turn}
\begin{itemize}
\raggedright
  \item \emph{Say it:} \emph{ʿamīq}
  \item \emph{Say it:} \emph{ʿamīq}, once more
  \item \emph{Say it:} say \emph{ʿamīq}
  \item \emph{Recall:} say the Arabic for pleased, then the Arabic for spacious, then say \emph{ʿamīq} again
\end{itemize}

\begin{tcolorbox}[breakable,colback=teal!4,colframe=teal!35!black,title={Before you move on}]
What does \ar{عميق} mean? (\textquotedblleft{}Deep\textquotedblright{}.) Say it once more.
\end{tcolorbox}
\section[\texorpdfstring{\ar{مستدير} (mustadīr)}{مستدير (mustadīr)}]{\ar{مستدير} (mustadīr) — round}

\label{lesson:AR-C148-mustadir}



\begin{quote}
Before the new one: say the Arabic for spacious, then the Arabic for deep.
\end{quote}

\subsection*{You'll want to know: \ar{مستدير}}
\textbf{\ar{مستدير}} — \emph{mustadīr} — \textquotedblleft{}round\textquotedblright{}.

\begin{cousinweb}
\textbf{\ar{مستدير}} (\emph{mustadīr}) is built on the native Arabic root \emph{d-w-r}, \textquotedblleft{}to turn round\textquotedblright{} (Lane). \emph{Mustadīr} is the Form X participle, \textquotedblleft{}turning round\textquotedblright{}, hence \textquotedblleft{}round\textquotedblright{}.
\end{cousinweb}

\begin{grammarlens}[title={What you've built}]
One more word for this chapter.
\end{grammarlens}

\subsection*{Your turn}
\begin{itemize}
\raggedright
  \item \emph{Say it:} \emph{mustadīr}
  \item \emph{Say it:} \emph{mustadīr}, once more
  \item \emph{Say it:} say \emph{mustadīr}
  \item \emph{Recall:} say the Arabic for spacious, then the Arabic for deep, then say \emph{mustadīr} again
\end{itemize}

\begin{tcolorbox}[breakable,colback=teal!4,colframe=teal!35!black,title={Before you move on}]
What does \ar{مستدير} mean? (\textquotedblleft{}Round\textquotedblright{}.) Say it once more.
\end{tcolorbox}
\section[\texorpdfstring{\ar{حاد} (ḥādd)}{حاد (ḥādd)}]{\ar{حاد} (ḥādd) — sharp}

\label{lesson:AR-C148-hadd}



\begin{quote}
Before the new one: say the Arabic for deep, then the Arabic for round.
\end{quote}

\subsection*{You'll want to know: \ar{حاد}}
\textbf{\ar{حاد}} — \emph{ḥādd} — \textquotedblleft{}sharp\textquotedblright{}.

\begin{cousinweb}
\textbf{\ar{حاد}} (\emph{ḥādd}) is built on the native Arabic root \emph{ḥ-d-d}, \textquotedblleft{}to sharpen; a limit\textquotedblright{} (Lane). \emph{Ḥādd} is the active participle, \textquotedblleft{}sharp\textquotedblright{}.
\end{cousinweb}

\begin{grammarlens}[title={What you've built}]
One more word for this chapter.
\end{grammarlens}

\subsection*{Your turn}
\begin{itemize}
\raggedright
  \item \emph{Say it:} \emph{ḥādd}
  \item \emph{Say it:} \emph{ḥādd}, once more
  \item \emph{Say it:} say \emph{ḥādd}
  \item \emph{Recall:} say the Arabic for deep, then the Arabic for round, then say \emph{ḥādd} again
\end{itemize}

\begin{tcolorbox}[breakable,colback=teal!4,colframe=teal!35!black,title={Before you move on}]
What does \ar{حاد} mean? (\textquotedblleft{}Sharp\textquotedblright{}.) Say it once more.
\end{tcolorbox}
